\chapter{SELINGAN SANGAT SINGKAT TENTANG BAHASA KATEGORI}\label{categories}

Dalam matematika kita mempelajari berbagai hal (objek) dan pemetaan tertentu di antara hal-hal tersebut (morfisme).
Beberapa contohnya ialah himpunan dan fungsi, grup dan homomorfisme, ruang topologis
dan pemetaan kontinu, ruang vektor dan transformasi linear, serta ruang Hilbert dan
pemetaan linear terbatas. Contoh-contoh ini masing-masing berasal dari cabang matematika yang berbeda---teori
himpunan, teori grup, topologi, aljabar linear, dan analisis fungsional.
Namun, bidang-bidang yang berbeda ini memiliki banyak kesamaan: istilah seperti hasil kali,
koproduk, subobjek, hasil bagi, tarik balik, isomorfisme, dan limit proyektif muncul dalam banyak bidang.
Teori kategori merupakan upaya untuk menyatukan dan memformalkan sebagian konsep umum tersebut.
Dalam arti tertentu, teori kategori mempelajari hal-hal yang sama-sama dimiliki oleh berbagai cabang matematika.
Barangkali uraian yang lebih baik ialah: bahasa kategori memberi tahu kita ``bagaimana sesuatu bekerja'',
bukan ``apa sesuatu itu''.

Catatan ini---dan memang setiap buku pada tingkat ini---menggunakan bahasa himpunan di mana-mana tanpa
mensyaratkan kajian terperinci tentang teori himpunan aksiomatik. Demikian pula, kita akan dengan leluasa menggunakan
\emph{bahasa} kategori tanpa terlebih dahulu memulai kajian teori kategori yang memuaskan dari segi
fondasi (teori kategori sendiri merupakan bidang penelitian yang luas dan penting). Kadang-kadang buku teks membuat
bahkan bahasa kategori sulit dipelajari karena sangat mengandalkan latar belakang aljabar pembaca.
Sebagaimana kebanyakan orang merasa nyaman menggunakan bahasa himpunan (terlepas dari apakah mereka pernah mengkaji
\emph{teori} himpunan secara serius), hampir setiap orang seharusnya mendapati bahwa bahasa kategori nyaman digunakan
dan memberi pencerahan tanpa prasyarat yang rumit.

Bagi pembaca yang ingin mendalami pokok ini, saya merekomendasikan pengantar yang sangat ramah ke dunia
kategori, yang terdapat sebagai Bab 3 dalam buku indah karya Semadeni~\cite{Semadeni:1971}. Salah satu buku klasik
yang ditulis oleh seorang pendiri bidang ini ialah~\cite{MacLane:1971}. Buku yang lebih baru ialah~\cite{LawvereS:1997}.

Dengan menunjukkan prinsip-prinsip pemersatu, bahasa kategori kerap memberikan wawasan mencolok tentang
``cara sesuatu bekerja'' dalam matematika. Yang sama pentingnya, kita menjadi lebih efisien karena tidak perlu
mengulangi argumen yang pada hakikatnya sama dalam konteks yang hanya sedikit berbeda.

Untuk sementara, kita hanya akan mendefinisikan ``objek'', ``morfisme'', dan ``funktor'', lalu memberikan beberapa contoh.









\section{Objek dan Morfisme}\label{C0151}
\begin{defn} Misalkan $\sfml A$ suatu kelas yang anggota-anggotanya kita sebut
 \index{objek (dari suatu kategori)}%
\df{objek}. Untuk setiap pasangan objek $(S,T)$, kita kaitkan suatu kelas $\mor ST$ yang
anggota-anggotanya kita sebut
 \index{morfisme!dari suatu kategori}%
\df{morfisme} (atau
 \index{panah}%
\df{panah}) dari $S$ ke $T$. Kita mengasumsikan bahwa $\mor ST$ dan $\mor UV$ saling lepas,
kecuali jika $S = U$ dan $T = V$.

Kita mengandaikan pula adanya operasi $\circ$ (disebut \df{komposisi}) yang
mengaitkan setiap $\alpha \in \mor ST$ dan setiap $\beta \in \mor TU$ dengan suatu morfisme
$\beta \circ \alpha \in \mor SU$ sedemikian rupa sehingga:
 \begin{enumerate}
  \item $\gamma \circ (\beta \circ \alpha) = (\gamma \circ
\beta) \circ \alpha$ apabila $\alpha \in \mor ST$, $\beta \in \mor TU$, dan $\gamma \in
\mor UV$;
  \item untuk setiap objek $S$ terdapat morfisme $I_S \in
\mor SS$ yang memenuhi $\alpha \circ I_S = \alpha$ apabila $\alpha \in \mor ST$, dan $I_S
\circ \beta = \beta$ apabila $\beta \in \mor RS$.
 \end{enumerate}
Dalam keadaan ini, kelas $\sfml A$ bersama keluarga morfisme yang terkait merupakan suatu
 \index{kategori}%
\df{kategori}.

Kita akan mencadangkan notasi
 \index{<arrowto@$S \to^\alpha T$ (morfisme dalam suatu kategori)}%
$S \to^\alpha T$ untuk keadaan ketika $S$ dan $T$ merupakan objek dalam suatu kategori dan
$\alpha$ adalah morfisme anggota $\mor ST$. Seperti halnya pada grup dan ruang vektor,
kita biasanya menghilangkan simbol komposisi $\circ$ dan menulis
  \index{<binopconcat@$\beta\alpha$ (notasi untuk komposisi morfisme)}%
  \index{konvensi!notasi untuk komposisi!morfisme}%
$\beta\alpha$ untuk $\beta \circ \alpha$.

Suatu kategori disebut
 \index{kecil secara lokal}%
 \index{kategori!kecil secara lokal}%
\df{kecil secara lokal} jika $\mor ST$ merupakan himpunan untuk setiap pasangan objek $S$ dan $T$ dalam kategori tersebut. Kategori itu disebut
 \index{kecil}%
 \index{kategori!kecil}%
\df{kecil} jika, selain itu, kelas semua objeknya merupakan suatu himpunan. Dalam catatan ini, kategori-kategori yang kita minati
akan bersifat kecil secara lokal.
\end{defn}

\begin{exam} Kategori
 \index{kategori!$\cat{SET}$ sebagai suatu}%
 \index{SET@$\cat{SET}$!kategori}%
$\cat{SET}$ memiliki himpunan sebagai objek dan fungsi (pemetaan) sebagai morfisme.
\end{exam}

\begin{exam} Kategori
 \index{kategori!$\cat{AbGp}$ sebagai suatu}%
 \index{AbGp@$\cat{AbGp}$!kategori}%
$\cat{AbGp}$ memiliki grup Abelian sebagai objek dan homomorfisme grup sebagai morfisme.
\end{exam}

\begin{exam} Kategori
 \index{kategori!$\cat{VEC}$ sebagai suatu}%
 \index{VEC@$\cat{VEC}$!kategori}%
$\cat{VEC}$ memiliki ruang vektor sebagai objek dan transformasi linear sebagai morfisme.
\end{exam}







\begin{defn}\label{C008111} Misalkan $\le$ suatu relasi pada himpunan tak kosong $S$.
 \begin{enumerate}
  \item Jika relasi
 \index{<le@$\le$ (notasi untuk suatu urutan parsial)}%
 $\le$ bersifat refleksif (yakni, $x \le x$ untuk setiap $x \in S$) dan transitif (yakni, jika $x \le y$ dan $y \le z$, maka $x  \le z$), relasi itu disebut suatu
 \index{refleksif}%
 \index{transitif}%
 \index{praurutan}%
 \index{urutan!pra-}%
\df{praurutan}.
  \item Jika $\le$ merupakan praurutan dan juga bersifat antisimetris (yakni, jika $x \le y$ dan $y \le x$, maka $x = y$), relasi itu disebut suatu
 \index{antisimetris}%
 \index{parsial!urutan}%
 \index{urutan!parsial}%
\df{urutan parsial}.
  \item Elemen $x$ dan $y$ dalam suatu himpunan yang telah diberi praurutan disebut
 \index{dapat dibandingkan}%
\df{dapat dibandingkan} jika $x \le y$ atau $y \le x$.
  \item Jika $\le$ merupakan urutan parsial yang membuat setiap dua elemen dapat dibandingkan, relasi itu disebut suatu
 \index{linear!urutan}%
 \index{urutan!linear}%
\df{urutan linear} (atau suatu
 \index{total!urutan}%
 \index{urutan!total}%
\df{urutan total}).
  \item Jika relasi $\le$ merupakan praurutan (berturut-turut, urutan parsial,
urutan linear) pada $S$, pasangan $(S,\le)$ disebut suatu
 \index{himpunan terpraurut}%
\df{himpunan terpraurut} (berturut-turut,
 \index{himpunan terurut parsial}%
\df{himpunan terurut parsial},
 \index{himpunan terurut linear}%
\df{himpunan terurut linear}).
  \item Subhimpunan terurut linear dari himpunan terurut parsial $(S,\le)$ disebut suatu
 \index{rantai}%
\df{rantai} di~$S$.
 \end{enumerate}
Kita boleh menulis $b \ge a$ sebagai pengganti $a \le b$. Notasi
 \index{<l@$<$ (notasi untuk suatu urutan ketat)}%
$a < b$ (atau, secara ekuivalen, $b > a$) berarti $a \le b$ dan $a \ne b$.
\end{defn}

\begin{defn} Pemetaan $f\colon S \sto T$ antara himpunan-himpunan terpraurut disebut
 \index{urutan!pemetaan pelestari}%
\df{pelestari urutan} jika $x \le y$ di $S$ mengakibatkan $f(x) \le f(y)$ di~$T$.
\end{defn}

\begin{exam} Kategori
 \index{kategori!$\cat{POSET}$ sebagai suatu}%
 \index{POSET@$\cat{POSET}$!kategori}%
$\cat{POSET}$ memiliki himpunan terurut parsial sebagai objek dan pemetaan pelestari urutan sebagai morfisme.
\end{exam}

\begin{notn} Misalkan $f \colon S \sto T$ suatu fungsi antara himpunan. Untuk $A \subseteq S$, kita mendefinisikan
$f^{\sto}(A) = \{f(x)\colon x \in A\}$, dan untuk $B \subseteq T$, kita mendefinisikan
 \index{<superscript@$f^{\sto}(A)$ atau $f(A)$ (citra $A$ di bawah $f$)}%
 \index{<superscript@$f^{\gets}(B)$ (pracitra $B$ di bawah $f$)}%
$f^{\gets}(B) = \{ x \in S\colon f(x) \in B\}$. Kita mengatakan bahwa $f^{\sto}(A)$ adalah
 \index{citra}%
\df{citra $A$ di bawah $f$} dan bahwa $f^{\gets}(B)$ adalah
 \index{pracitra}%
\df{pracitra} (atau
 \index{invers!citra}%
 \index{citra!invers}%
\df{citra invers}) \df{$B$ di bawah~$f$}. Kadang-kadang, untuk menghindari notasi yang berjejal, saya akan menyerah pada kebiasaan meragukan dengan
menulis $f(A)$ sebagai pengganti~$f^{\sto}(A)$.
\end{notn}

\begin{defn} Ingat bahwa keluarga $\sfml T$ yang terdiri atas subhimpunan-subhimpunan suatu himpunan $X$ disebut suatu
 \index{topologi}%
\df{topologi} pada $X$ jika keluarga itu memuat himpunan kosong $\emptyset$ dan $X$ sendiri serta tertutup terhadap pembentukan gabungan
($\bigcup\sfml S \in \sfml T$ apabila $\sfml S \subseteq \sfml T$) dan irisan hingga ($\bigcap\sfml F \in \sfml T$
apabila $\sfml F \subseteq \sfml T$ dan $\sfml F$ berhingga). Himpunan tak kosong $X$ bersama suatu topologi pada $X$ disebut suatu
 \index{topologis!ruang}%
 \index{ruang!topologis}%
\df{ruang topologis}. Jika $(X,\sfml T)$ suatu ruang topologis, setiap anggota $\sfml T$ disebut suatu
 \index{terbuka}%
 \index{<binrel@$\open UX$ ($U$ merupakan subhimpunan terbuka dari~$X$)}%
\df{himpunan terbuka}. Untuk menyatakan bahwa himpunan $U$ merupakan subhimpunan terbuka dari $X$, kita sering menulis $\open UX$. Komplemen dari suatu himpunan terbuka disebut suatu
 \index{tertutup}%
\df{himpunan tertutup}. Fungsi $f\colon X \sto Y$ antara ruang-ruang topologis disebut
 \index{kontinu!fungsi}%
 \index{fungsi!kontinu}%
\df{kontinu} jika pracitra setiap himpunan terbuka bersifat terbuka; yakni, jika $\open{f^{\gets}(V)}X$ apabila $\open VY$.
(Jika Anda belum pernah menjumpai ruang topologis, lihat dua butir pertama dalam bagian pertama
Bab 10 dan bagian pertama Bab 11 pada catatan daring saya~\cite{Erdman:2007}.)
\end{defn}

\begin{exam} Setiap himpunan tak kosong $X$ memiliki topologi terkecil (terkasar); topologi itu terdiri tepat atas dua himpunan,
$\emptyset$ dan~$X$. Topologi ini adalah topologi
 \index{indiskret}%
 \index{topologi!indiskret}%
\df{indiskret} pada~$X$. (Perhatikan ejaan istilah bahasa Inggrisnya: terdapat kata lain dalam bahasa Inggris, ``indiscreet'',
yang bermakna sama sekali berbeda.) Himpunan $X$ juga memiliki topologi terbesar (terhalus), yang terdiri atas
keluarga $\sfml P(X)$ dari semua subhimpunan~$X$. Topologi ini adalah topologi
 \index{diskret}%
 \index{topologi!diskret}%
\df{diskret} pada~$X$.
  \end{exam}

\begin{exam} Kategori
 \index{kategori!$\cat{TOP}$ sebagai suatu}%
 \index{TOP@$\cat{TOP}$!kategori}%
$\cat{TOP}$ memiliki ruang topologis sebagai objek dan fungsi kontinu sebagai morfisme.
\end{exam}

\begin{defn}\label{defn_alg} Misalkan $(A,+,M)$ suatu ruang vektor atas $\K$ yang dilengkapi dengan
operasi biner lain $A \times A \sto A$, dengan $(x,y) \mapsto x \cdot y$, sedemikian
rupa sehingga $(A,+,\cdot\,)$ merupakan gelanggang. (Notasi $x \cdot y$ biasanya disingkat menjadi $xy$.) Jika, selain itu, persamaan
  \begin{equation}\label{eq_def_algebra}
     \alpha (xy) = (\alpha x)y = x(\alpha y)
  \end{equation}
berlaku untuk semua $x$, $y \in A$ dan $\alpha \in \K$, maka $(A,+,M,\,\cdot\,)$ disebut suatu
 \index{aljabar}%
\df{aljabar} atas medan~$\K$ (kadang-kadang disebut \df{aljabar asosiatif
linear}). Seperti biasa, kita menyalahgunakan terminologi dengan menulis hal-hal seperti, ``Misalkan $A$ suatu
aljabar.'' Kita mengatakan bahwa suatu aljabar $A$
 \index{beridentitas!aljabar}%
 \index{aljabar!beridentitas}%
\df{beridentitas} jika gelanggang yang mendasarinya $(A,+,\cdot\,)$ memiliki identitas perkalian; yakni, jika terdapat
 \index{perkalian!identitas}%
 \index{identitas!perkalian}%
suatu elemen $\vc 1_A \ne \vc 0$ di $A$ sehingga $\vc 1_A \cdot x = x \cdot \vc 1_A = x$ untuk semua $x \in A$.
  Aljabar itu disebut
 \index{komutatif!aljabar}%
 \index{aljabar!komutatif}%
\df{komutatif} jika gelanggangnya komutatif; yakni, jika $xy = yx$ untuk semua $x$, $y \in A$.

Subhimpunan $B$ dari suatu aljabar $A$ disebut suatu
 \index{subaljabar}%
\df{subaljabar} dari $A$ jika subhimpunan tersebut merupakan aljabar dengan operasi-operasi yang diwarisinya dari~$A$. Suatu
subaljabar $B$ dari aljabar beridentitas $A$ disebut
 \index{beridentitas!subaljabar}%
 \index{subaljabar!beridentitas}%
\df{subaljabar beridentitas} jika memuat identitas perkalian dari~$A$.

\textbf{PERHATIAN.} Agar menjadi subaljabar beridentitas, 
\emph{tidaklah} cukup jika $B$ memiliki identitas perkaliannya sendiri; subaljabar itu harus memuat identitas dari~$A$. Jadi, suatu
aljabar dapat sekaligus beridentitas dan merupakan subaljabar dari $A$ tanpa menjadi subaljabar beridentitas
dari~$A$.

Pemetaan $f\colon A \sto B$ antara aljabar-aljabar disebut suatu
 \index{homomorfisme!aljabar}%
 \index{aljabar!homomorfisme}%
\df{homomorfisme (aljabar)} jika pemetaan itu merupakan pemetaan linear antara $A$ dan $B$ sebagai ruang vektor yang mempertahankan perkalian
(yakni, $f(xy) = f(x)f(y)$ untuk semua $x$, $y \in A$). Dengan kata lain, homomorfisme aljabar adalah homomorfisme gelanggang linear.
Pemetaan itu disebut
 \index{beridentitas!homomorfisme aljabar}%
 \index{aljabar!homomorfisme!beridentitas}%
 \index{homomorfisme!aljabar!beridentitas}%
\df{homomorfisme (aljabar) beridentitas} jika mempertahankan identitas; yakni, jika $A$ dan $B$ keduanya merupakan aljabar beridentitas dan $f(\vc 1_A) = \vc 1_B$.
 \index{kernel!homomorfisme aljabar}%
\df{Kernel} dari suatu homomorfisme aljabar $f\colon A \sto B$ tentu saja adalah $\{a \in A \colon f(a) = \vc 0\}$.

Jika $f^{-1}$ ada dan juga merupakan homomorfisme aljabar, maka $f$ disebut suatu
 \index{isomorfisme!aljabar}%
\df{isomorfisme} dari $A$ ke $B$. Jika terdapat isomorfisme dari $A$ ke $B$, maka $A$ dan
$B$ disebut
 \index{isomorfik}%
\df{isomorfik}.
\end{defn}

\begin{exam} Kategori
 \index{kategori!$\cat{ALG}$ sebagai suatu}%
 \index{ALG@$\cat{ALG}$!kategori}%
$\cat{ALG}$ memiliki aljabar sebagai objek dan homomorfisme aljabar sebagai morfisme.
\end{exam}

Contoh-contoh sebelumnya merupakan contoh \emph{kategori konkret}---yakni, kategori
yang objek-objeknya berupa himpunan (biasanya bersama struktur tambahan) dan
morfisme-morfismenya berupa fungsi (biasanya dengan mempertahankan struktur tambahan tersebut). Dalam catatan ini,
kategori-kategori yang kita minati bersifat konkret sekaligus kecil secara lokal. Berikut ini (bagi pembaca yang ingin tahu) adalah
definisi yang lebih formal untuk \emph{kategori konkret}.

\begin{defn}\label{C015124} Suatu kategori $\sfml A$ bersama fungsi $\abs{\hphantom{O}}$ yang mengaitkan
setiap objek $A$ di $\sfml A$ dengan suatu himpunan $\abs A$ disebut suatu
 \index{kategori konkret}%
 \index{kategori!konkret}%
\df{kategori konkret} jika syarat-syarat berikut dipenuhi:
 \begin{enumerate}
    \item setiap morfisme $A \to^f B$ merupakan fungsi dari $\abs A$ ke~$\abs B$;
    \item setiap morfisme identitas $I_A$ di $\sfml A$ merupakan fungsi identitas pada~$\abs A$; dan
    \item komposisi morfisme $A \to^f B$ dan $B \to^g C$ sama dengan komposisi
fungsi $f\colon \abs A \sto \abs B$ dan $g\colon \abs B \sto \abs C$.
 \end{enumerate}
Jika $A$ suatu objek dalam kategori konkret $\sfml A$, maka $\abs A$ disebut
 \index{himpunan dasar}%
 \index{himpunan!dasar}%
 \index{<unop@$\abs A$ (himpunan dasar dari objek~$A$)}%
\df{himpunan dasar} dari~$A$.
\end{defn}

Meskipun kategori-kategori yang kita minati dalam catatan ini memang merupakan kategori
konkret, mungkin tetap menarik untuk melihat contoh kategori yang
\emph{tidak disajikan secara konkret}, yakni tanpa terlebih dahulu mewujudkan objek sebagai himpunan dan morfisme sebagai fungsi.

\begin{exam}\label{C015127} Misalkan $G$ suatu monoid (yakni, semigrup beridentitas).
 \index{monoid}%
Tinjau suatu kategori $\cat C$ yang memiliki tepat satu objek, yang kita sebut~$\star$. Karena hanya
ada satu objek, hanya ada satu keluarga morfisme $\mor{\star}{\star}$, yang kita
ambil sebagai~$G$. Komposisi morfisme didefinisikan sebagai perkalian monoid.
Artinya, $a \circ b := ab$ untuk semua $a$, $b \in G$. Jelas bahwa komposisi bersifat asosiatif dan
elemen identitas $G$ merupakan morfisme identitas. Jadi, $\cat C$ adalah suatu kategori.
\end{exam}

\begin{defn} Dalam setiap kategori konkret, kita akan menyebut morfisme injektif sebagai suatu
 \index{monomorfisme}%
\df{monomorfisme} dan morfisme surjektif sebagai suatu
 \index{epimorfisme}%
\df{epimorfisme}.
\end{defn}

\begin{cau} Definisi-definisi di atas mencerminkan penggunaan asli istilah tersebut oleh Bourbaki dan merupakan definisi yang
paling lazim dipakai oleh matematikawan di luar teori kategori itu sendiri; di dalam teori kategori,
``monomorfisme'' berarti ``dapat dibatalkan dari kiri'' dan ``epimorfisme'' berarti ``dapat
dibatalkan dari kanan''. (Perhatikan bahwa istilah \emph{injektif} dan \emph{surjektif} mungkin tidak
bermakna jika diterapkan pada morfisme dalam kategori yang tidak konkret.)

Suatu morfisme $B \to^g C$ disebut
 \index{kiri!dapat dibatalkan dari}%
 \index{dapat dibatalkan!dari kiri}%
\df{dapat dibatalkan dari kiri} jika, setiap kali morfisme $A \to^{f_1} B$ dan $A \to^{f_2} B$ memenuhi
$gf_1 = gf_2$, berlaku $f_1 = f_2$. Mac Lane menyarankan agar morfisme yang dapat dibatalkan dari kiri disebut
 \index{monik}%
morfisme \df{monik}. Perbedaan antara morfisme monik dan monomorfisme ternyata
kecil. Dalam catatan ini, hampir semua morfisme yang kita jumpai bersifat monik jika dan
hanya jika merupakan monomorfisme. Sebagai latihan mudah, buktikan bahwa setiap morfisme injektif
dalam suatu kategori (konkret) bersifat monik. Konversnya kadang-kadang gagal.

Dengan semangat yang sama, Mac Lane menyarankan agar suatu morfisme yang
 \index{kanan!dapat dibatalkan dari}%
 \index{dapat dibatalkan!dari kanan}%
\emph{dapat dibatalkan dari kanan} (yakni, suatu morfisme $A \to^f B$ sedemikian sehingga, apabila
morfisme $B \to^{g_1} C$ dan $B \to^{g_2} C$ memenuhi $g_1f = g_2f$, maka $g_1 = g_2$) disebut suatu
 \index{epik}%
morfisme \df{epik}. Sekali lagi, merupakan latihan mudah untuk menunjukkan bahwa dalam suatu kategori (konkret),
setiap epimorfisme bersifat epik. Namun, konversnya gagal dalam beberapa kategori yang cukup umum.
\end{cau}

\begin{defn} Terminologi untuk invers morfisme dalam kategori pada dasarnya sama
dengan terminologi untuk fungsi. Misalkan $A \to^\alpha B$ dan $B \to^\beta A$ merupakan
morfisme dalam suatu kategori. Jika $\beta \circ \alpha = I_A$, maka $\beta$ adalah suatu
 \index{kiri!invers!morfisme}%
 \index{invers!kiri}%
\df{invers kiri} dari $\alpha$ dan, secara ekuivalen, $\alpha$ adalah suatu
 \index{kanan!invers!morfisme}%
 \index{invers!kanan}%
\df{invers kanan} dari $\beta$. Kita mengatakan bahwa morfisme $\alpha$ merupakan suatu
 \index{isomorfisme!dalam suatu kategori}%
\df{isomorfisme} (atau bersifat
 \index{invertibel!morfisme}%
\df{invertibel}) jika terdapat morfisme $B \to^\beta A$ yang sekaligus merupakan invers kiri dan
invers kanan dari~$\alpha$. Morfisme demikian dinotasikan dengan $\alpha^{-1}$ dan disebut
 \index{<unopur@$\alpha^{-1}$ (invers morfisme $\alpha$)}%
 \index{invers!morfisme}%
\df{invers} dari~$\alpha$.
\end{defn}

\begin{prop} Misalkan $A \to^\alpha B$ suatu morfisme dalam suatu kategori. Jika $\alpha$ memiliki invers kiri $\lambda$ dan
invers kanan $\rho$, maka $\lambda = \rho$ dan, akibatnya, $\alpha$ merupakan suatu isomorfisme.
\end{prop}

Dalam setiap kategori konkret, kita dapat menanyakan apakah setiap morfisme bijektif (yakni, setiap
pemetaan yang sekaligus merupakan monomorfisme dan epimorfisme) adalah suatu isomorfisme. Jawabannya
sering kali merupakan \emph{ya} yang mudah (seperti dalam $\cat{SET}$, $\cat{AbGp}$, dan $\cat{VEC}$) atau
\emph{tidak} yang mudah (misalnya, dalam kategori $\cat{POSET}$
 \index{morfisme bijektif!tidak harus invertibel!dalam $\cat{POSET}$}%
 \index{POSET@$\cat{POSET}$!morfisme bijektif dalam}%
 \index{POSET@$\cat{POSET}$!kategori}%
yang terdiri atas himpunan terurut parsial dan pemetaan pelestari urutan). Namun,
sesekali jawabannya ternyata berupa hasil yang menarik dan mendalam (lihat, misalnya, \emph{teorema pemetaan
terbuka}~\futurexref{6.3.4}{C069414}).

\begin{exam} Dalam kategori $\cat{SET}$,
 \index{morfisme bijektif!invertibel!dalam $\cat{SET}$}%
 \index{SET@$\cat{SET}$!morfisme bijektif dalam}%
setiap morfisme bijektif merupakan suatu isomorfisme.
\end{exam}

\begin{exam}\label{C015216} Kategori
 \index{LAT@$\cat{LAT}$!kategori}%
 \index{kategori!$\cat{LAT}$ sebagai suatu}%
$\cat{LAT}$ memiliki kekisi sebagai objek dan homomorfisme
kekisi sebagai morfisme. (Suatu
 \index{kekisi}%
\df{kekisi} adalah himpunan terurut parsial yang setiap pasangan elemennya memiliki infimum dan supremum, dan suatu
 \index{kekisi!homomorfisme}%
 \index{homomorfisme!kekisi}%
\df{homomorfisme kekisi} adalah pemetaan antara kekisi-kekisi yang mempertahankan infimum dan supremum.)
Dalam kategori ini, morfisme bijektif
 \index{morfisme bijektif!invertibel!dalam $\cat{LAT}$}%
 \index{LAT@$\cat{LAT}$!morfisme bijektif dalam}%
merupakan isomorfisme.
\end{exam}

\begin{exam} Suatu
 \index{homeomorfisme}%
\df{homeomorfisme} adalah bijeksi kontinu $f$ antara ruang-ruang topologis yang inversnya $f^{-1}$ juga kontinu. Homeomorfisme
merupakan isomorfisme dalam kategori $\cat{TOP}$. Dalam kategori ini, tidak setiap bijeksi kontinu merupakan homeomorfisme.
\end{exam}

\begin{exam} Jika, dalam kategori $\cat C$ pada Contoh~\ref{C015127}, monoid $G$ merupakan suatu grup,
maka setiap morfisme dalam $\cat C$ merupakan suatu isomorfisme.
\end{exam}






\section{Funktor}
\begin{defn} Jika $\cat{A}$ dan $\cat{B}$ merupakan kategori, suatu
 \index{kovarian!funktor}%
 \index{funktor!kovarian}%
\df{funktor kovarian} $F$ dari $\cat A$ ke $\cat B$ (ditulis ${\cat A}\to^F{\cat B}$) adalah
sepasang pemetaan: suatu
 \index{objek!pemetaan}%
 \index{pemetaan!objek}%
\df{pemetaan objek} $F$ yang mengaitkan setiap objek $S$ dalam $\cat A$ dengan suatu objek $F(S)$ dalam
$\cat B$, dan suatu
 \index{morfisme!pemetaan}%
 \index{pemetaan!morfisme}%
\df{pemetaan morfisme} (yang juga dinotasikan dengan $F$) yang mengaitkan setiap morfisme $f \in \mor
ST$ dalam $\cat A$ dengan suatu morfisme $F(f) \in \mor {F(S)}{F(T)}$ dalam $\cat B$, sedemikian rupa sehingga
 \begin{enumerate}
  \item $F(g \circ f) = F(g) \circ F(f)$ apabila $g \circ f$ terdefinisi dalam $\cat A$; dan
  \item $F(\id S) = \id{F(S)}$ untuk setiap objek $S$ dalam $\cat A$.
 \end{enumerate}

Definisi suatu
 \index{kontravarian!funktor}%
 \index{funktor!kontravarian}%
\df{funktor kontravarian} ${\cat A}\to^F{\cat B}$ berbeda dari definisi sebelumnya
hanya dalam dua hal. Pertama, pemetaan morfismenya mengaitkan setiap morfisme $f \in \mor ST$ dalam
$\cat A$ dengan suatu morfisme $F(f) \in \mor {F(T)}{F(S)}$ dalam $\cat B$. Kedua, syarat (a)
di atas diganti dengan
 \begin{enumerate}
  \item[(\rm{a'})] $F(g \circ f) = F(f) \circ F(g)$ apabila $g \circ f$
terdefinisi dalam~$\cat A$.
 \end{enumerate}
\end{defn}

\begin{exam}\label{functors005exam} Suatu
 \index{pelupa!funktor}%
 \index{funktor!pelupa}%
\df{funktor pelupa} adalah funktor yang memetakan objek dan morfisme dari suatu kategori $\cat
C$ ke kategori $\cat C'$ yang memiliki struktur atau sifat lebih sedikit. Sebagai contoh, jika $V$
suatu ruang vektor, funktor $F$ yang ``melupakan'' operasi perkalian skalar
pada ruang vektor akan memetakan $V$ ke kategori grup Abelian. (Grup
Abelian $F(V)$ akan memiliki himpunan elemen yang sama dengan ruang vektor $V$ dan
operasi penjumlahan yang sama, tetapi tidak memiliki perkalian skalar.) Suatu pemetaan linear
$T\colon V \sto W$ antara ruang-ruang vektor akan dibawa oleh funktor $F$ menjadi homomorfisme
grup $F(T)$ antara grup-grup Abelian $F(V)$ dan~$F(W)$.

Funktor pelupa juga dapat ``melupakan'' sifat. Jika $G$ suatu objek dalam
kategori grup Abelian, funktor yang ``melupakan'' komutativitas grup Abelian
akan membawa $G$ ke kategori grup.

Pada bagian sebelumnya telah disebutkan bahwa semua kategori yang kita minati dalam
catatan ini merupakan kategori konkret (kategori yang objek-objeknya berupa himpunan dengan struktur
tambahan dan morfisme-morfismenya berupa pemetaan yang, dalam pengertian tertentu, mempertahankan
struktur tambahan tersebut). Kita akan beberapa kali menggunakan jenis khusus funktor
pelupa---funktor yang melupakan seluruh struktur objek kecuali himpunan dasarnya dan yang
melupakan segala sifat morfisme yang mempertahankan struktur. Jika $A$ suatu
objek dalam suatu kategori konkret $\cat C$, kita menotasikan
 \index{<unop@$\abs A$ (funktor pelupa yang bekerja pada~$A$)}%
himpunan dasarnya dengan $\abs{A}$. Jika \mbox{$A~\to^f~B$} suatu morfisme dalam $\cat C$, kita
menotasikan dengan $\abs f$ pemetaan dari $\abs A$ ke $\abs B$ yang dipandang sekadar sebagai fungsi
antara himpunan. Mudah dilihat bahwa $\abs{\hphantom{M}}$\,, yang membawa objek dalam
$\cat C$ ke objek dalam $\cat{SET}$ (kategori himpunan dan pemetaan) serta morfisme dalam $\cat
C$ ke morfisme dalam $\cat{SET}$, merupakan suatu funktor kovarian.

Dalam kategori $\cat{VEC}$ yang terdiri atas ruang vektor dan pemetaan linear, misalnya,
$\abs{\hphantom{M}}$ membuat ruang vektor $V$ ``melupakan'' penjumlahan maupun
perkalian skalarnya ($\abs V$ hanyalah suatu himpunan). Jika $T\colon V \sto W$ suatu transformasi
linear, maka $\abs T\colon \abs V \sto \abs W$ hanyalah pemetaan antara himpunan---pemetaan itu telah
``melupakan'' sifat mempertahankan operasi.
\end{exam}


\begin{defn} Suatu himpunan terurut parsial disebut
 \index{urutan!lengkap}%
 \index{lengkap!urutan}%
\df{lengkap secara urutan} jika setiap subhimpunan tak kosong memiliki supremum (yakni, batas atas
terkecil) dan infimum (batas bawah terbesar).
\end{defn}

\begin{defn} Misalkan $S$ suatu himpunan. Maka
 \index{himpunan kuasa}%
 \index{P@$\sfml P(S)$ (himpunan kuasa dari $S$)}%
\df{himpunan kuasa} dari $S$, yang dinotasikan dengan $\sfml P(S)$, adalah keluarga semua subhimpunan~$S$.
\end{defn}

\begin{exam}[Funktor himpunan kuasa] Misalkan $S$ suatu himpunan tak kosong.
 \begin{enumerate}
  \item Himpunan kuasa $\sfml P(S)$ dari $S$ yang diurutkan secara parsial oleh $\subseteq$
bersifat lengkap secara urutan.
  \item Kelas himpunan terurut parsial yang lengkap secara urutan beserta pemetaan
pelestari urutan merupakan suatu kategori.
 \index{funktor!himpunan kuasa}%
 \index{himpunan kuasa!funktor}%
  \item Untuk setiap fungsi $f$ antara himpunan, misalkan $\sfml P(f) = f^{\sto}$. Maka
$\sfml P$ merupakan funktor kovarian dari kategori himpunan dan fungsi ke kategori
himpunan terurut parsial yang lengkap secara urutan beserta pemetaan pelestari urutan.
  \item Untuk setiap fungsi $f$ antara himpunan, misalkan $\sfml P(f) = f^{\gets}$. Maka
$\sfml P$ merupakan funktor kontravarian dari kategori himpunan dan fungsi ke kategori
himpunan terurut parsial yang lengkap secara urutan beserta pemetaan pelestari urutan.
 \end{enumerate}
\end{exam}

\begin{defn}\label{defn_vsp_adjoint} Misalkan $T\colon V \sto W$ suatu pemetaan linear antara ruang-ruang vektor. Untuk setiap
$g \in W^{\#}$, misalkan $T^{\#}(g) = g\,T$. Perhatikan bahwa $T^{\#}(g) \in V^{\#}$. Pemetaan $T^{\#}$ dari ruang
vektor $W^{\#}$ ke ruang vektor~$V^{\#}$ disebut pemetaan
 \index{adjoin!ruang vektor}%
 \index{vektor!ruang!pemetaan adjoin}%
 \index{<unopur@$T^\#$ (adjoin ruang vektor)}%
\df{adjoin} (ruang vektor) dari~$T$.
\end{defn}

\begin{exam} Pasangan pemetaan $V \mapsto V^{\#}$ (yang membawa suatu ruang vektor ke dual aljabarnya) dan
$T \mapsto T^{\#}$ (yang membawa suatu pemetaan linear ke dualnya) merupakan funktor kontravarian dari kategori $\cat{VEC}$
yang terdiri atas ruang vektor dan pemetaan linear ke kategori itu sendiri.
\end{exam}

\begin{exam}\label{0007606} Jika $\cat C$ suatu kategori, definisikan
 \index{kategori@$\cat{C^2}$ (pasangan objek dan morfisme)}%
$\cat{C^2}$ sebagai kategori yang objek-objeknya merupakan pasangan terurut objek dalam $\cat C$ dan
morfisme-morfismenya merupakan pasangan terurut morfisme dalam~$\cat C$. Jadi, jika $A \to^f C$ dan $B
\to^g D$ merupakan morfisme dalam $\cat C$, maka $(A,B) \to^{(f,g)} (C,D)$ (dengan $(f,g)(a,b) =
\bigl(f(a),g(b)\bigr)$ untuk semua $a \in A$ dan $b \in B$ apabila kategori tersebut konkret) merupakan morfisme dalam~$\cat{C^2}$.
Komposisi morfisme didefinisikan dengan cara yang wajar: $(f,g) \circ (h,j) = (f \circ h, g
\circ j)$. Kita mendefinisikan
 \index{diagonal!funktor}%
 \index{funktor!diagonal}%
\df{funktor diagonal} $\cat C \to^{\ftr D} \cat{C^2}$ dengan $\ftr D(A) := (A,A)$ dan
$\ftr D(f) := (f,f)$. Ini merupakan
funktor kovarian.
\end{exam}

\begin{exam}\label{C029717} Misalkan $X$ dan $Y$ ruang-ruang topologis dan $\phi\colon X \sto Y$
kontinu. Definisikan $\fml C \phi$ pada $\fml C(Y)$ dengan
    \[ \fml C \phi(g) = g \circ \phi \]
untuk semua $g \in \fml C(Y)$. Maka pasangan pemetaan $X \mapsto \fml C(X)$ dan
 \index{C@$\fml C$ (funktor)}%
$\phi \mapsto \fml C \phi$ merupakan funktor kontravarian dari kategori $\cat{TOP}$ yang terdiri atas ruang
topologis dan pemetaan kontinu ke kategori aljabar beridentitas dan homomorfisme
aljabar beridentitas.
\end{exam}











\endinput
